\subsubsection{Data Processing}
\begin{enumerate}
    \item \textbf{Normalization of Expression Values}
          
          The function runs a complete normalization pipeline on gene expression data.
          It first applies standard deviation normalization, then quantile normalization,
          followed by replicate averaging, and finally a log2 transformation. The
          result is returned in an array, which contains the fully normalized expression
          values. Call the \texttt{tox\_normalization\_pipeline} function with the following
          arguments:
          
          \begin{itemize}
            \item \texttt{input\_matrix}: Numeric expression vector matrix with size
                  \texttt{genes $\times$ tissues}
                  
            \item \texttt{group\_s}: Integer vector: start column index for each
                  replicate group (1-based)
                  
            \item \texttt{group\_c}: Integer vector: number of columns per replicate
                  group
          \end{itemize}
          
          \textbf{Return:}
          \begin{itemize}
            \item Numeric matrix: log2(x+1) normalized expression
          \end{itemize}
          
          \begin{lstlisting}[language=R, caption={R Normalization Pipeline}]
tox_normalization_pipeline(input_matrix, group_s, group_c)
          \end{lstlisting}
          
          \textbf{Alternatively you can calculate each step of the normalization
          pipeline manually by calling the following functions in sequence:}
          \begin{enumerate}
            \item \textbf{Normalize by Standard Deviation}
                  
                  The function normalizes each gene's expression vector using \texttt{sqrt(mean(x\textasciicircum
                    2))} across tissues (not classical standard deviation). Call the \texttt{tox\_normalize\_by\_std\_dev}
                  function with the following arguments:
                  
                  \begin{itemize}
                    \item \texttt{input\_matrix}: A numeric matrix with genes as rows and
                          tissues as columns
                  \end{itemize}
                  
                  \textbf{Returns:}
                  \begin{itemize}
                    \item A normalized numeric matrix with the same dimensions and names
                          as the input
                  \end{itemize}
                  \begin{lstlisting}[language=R, caption={R Standard Deviation Normalization}]
tox_normalize_by_std_dev(input_matrix)
                  \end{lstlisting}
                  
            \item \textbf{Quantile Normalization}
                  
                  The function performs quantile normalization of a gene expression matrix
                  (F42-compliant) and computes average expression per rank across tissues.
                  Call the \texttt{tox\_quantile\_normalization} function with the following
                  arguments:
                  
                  \begin{itemize}
                    \item \texttt{input\_matrix}: A numeric matrix with genes as rows and
                          tissues as columns
                  \end{itemize}
                  
                  \textbf{Returns:}
                  \begin{itemize}
                    \item A quantile-normalized numeric matrix with the same dimensions and
                          names as the input
                  \end{itemize}
                  \begin{lstlisting}[language=R, caption={R Quantile Normalization}]
tox_quantile_normalization(input_matrix)
                  \end{lstlisting}
                  
            \item \textbf{Calculate Tissue Averages}
                  
                  The function calculates tissue averages by averaging replicates within
                  each group. For each group of tissue replicates, it computes the average
                  expression per gene. The input matrix is column-major, flattened as a 1D
                  array. Call the \texttt{tox\_calculate\_tissue\_averages} function with
                  the following arguments:
                  
                  \begin{itemize}
                    \item \texttt{df}: A data frame or matrix with genes as rows and tissue
                          replicates as columns
                  \end{itemize}
                  
                  \textbf{Returns:}
                  \begin{itemize}
                    \item A data frame with genes as rows and averaged tissues as columns
                  \end{itemize}
                  \begin{lstlisting}[language=R, caption={R Tissue Averages Calculation}]
tox_calculate_tissue_averages(df)
                  \end{lstlisting}
                  
            \item \textbf{Log2 Transformation}
                  
                  The function applies \texttt{log2(x + 1)} transformation to each
                  element of the input matrix. This is computed via \texttt{log(x + 1) /
                    log(2)} for compatibility with WebAssembly (WASM). Call the \texttt{tox\_log2\_transformation}
                  function with the following arguments:
                  
                  \begin{itemize}
                    \item \texttt{input\_matrix}: A numeric matrix with genes as rows and
                          tissues as columns
                  \end{itemize}
                  
                  \textbf{Returns:}
                  \begin{itemize}
                    \item A numeric matrix with log2-transformed expression values,
                          preserving the same dimensions and names as the input.
                  \end{itemize}
                  \begin{lstlisting}[language=R, caption={R Log2 Transformation}]
tox_log2_transformation(input_matrix)
                  \end{lstlisting}
          \end{enumerate}
          
          \textbf{Additionally there are two helpfer functions to improve the data
          quality and prevent calculation errors:}
          \begin{enumerate}
            \item \textbf{Diagnose Data Quality}
                  
                  The function checks your gene expression matrix for common data quality
                  issues, such as missing values (NA), infinite values, NaNs, zeros, and
                  negative values. It provides a summary of these problems, lists
                  affected genes, and reports basic statistics like minimum, maximum, and
                  mean values. Use this function to quickly assess the quality of your
                  data before running normalization or analysis steps. Call the \texttt{tox\_diagnose\_data\_quality}
                  function with the following arguments:
                  
                  \begin{itemize}
                    \item \texttt{input\_matrix}: A numeric matrix with genes as rows and
                          tissues as columns
                          
                    \item (Optional) \texttt{show\_details}: Logical indicating whether to
                          show detailed information (default is TRUE)
                  \end{itemize}
                  
                  \textbf{Returns:}
                  \begin{itemize}
                    \item A list with diagnostic information about the data quality
                  \end{itemize}
                  
                  \begin{lstlisting}[language=R, caption={R Diagnose Data Quality}]
tox_diagnose_data_quality(input_matrix, show_details)
                  \end{lstlisting}
                  
            \item \textbf{Clean Data for Normalization}
                  
                  The function prepares your gene expression matrix for normalization by
                  handling problematic values. It can remove or impute NA, NaN, and Inf
                  values, filter out genes with all zero values, and optionally convert
                  very small values to zero. You can choose different strategies for dealing
                  with missing data, ensuring your matrix is clean and ready for
                  downstream analysis. Call the \texttt{tox\_clean\_data\_for\_normalization}
                  function with the following arguments:
                  
                  \begin{itemize}
                    \item \texttt{df\_matrix}: A numeric matrix with genes as rows and tissues
                          as columns
                          
                    \item (Optional) \texttt{remove\_all\_zero\_genes}: Logical, whether
                          to remove genes that are all zeros (default = TRUE)
                          
                    \item (Optional) \texttt{na\_strategy}: Strategy for handling NA values:
                          "remove\_genes", "remove\_samples", "impute\_zero", "impute\_mean"
                          (default = "remove\_genes"
                          
                    \item (Optional) \texttt{min\_expression\_threshold}: Minimum expression
                          value to consider (values below this become 0 and default is 0)
                          
                    \item (Optional) \texttt{convert\_small\_to\_zero}: Logical, wheter
                          very small numbers should be converted to zero (default = FALSE)
                  \end{itemize}
                  
                  \textbf{Returns:}
                  \begin{itemize}
                    \item A cleaned matrix ready for normalization
                  \end{itemize}
                  
                  \begin{lstlisting}[language=R, caption={R Clean Data for Normalization}]
tox_clean_data_for_normalization(df_matrix, remove_all_zero_genes, 
    na_strategy, min_expression_threshold, convert_small_to_zero)
                  \end{lstlisting}
          \end{enumerate}
          
    \item \textbf{Calculate Fold Change}
          
          If fold change calculation is needed, this function can be used after
          normalization. The function calculates \texttt{log2} fold changes between
          condition and control columns. For each control-condition pair, this
          function computes the \texttt{log2} fold change by subtracting the expression
          value in the control column from the corresponding value in the condition column,
          for all genes. Call the \texttt{tox\_calculate\_fold\_changes} function
          with the following arguments:
          
          \begin{itemize}
            \item \texttt{df}: A data frame with genes as rows and tissues/conditions
                  as columns
                  
            \item \texttt{control\_pattern}: A string pattern to detect control columns
                  
            \item \texttt{condition\_patterns}: A character vector with patterns to detect
                  condition columns
          \end{itemize}
          
          \textbf{Returns:}
          \begin{itemize}
            \item A data frame with genes as rows and log2 fold change values as columns
          \end{itemize}
          
          \begin{lstlisting}[language=R, caption={R Fold Change Calculation}]
tox_calculate_fc_by_patterns(df, control_pattern, condition_patterns)
          \end{lstlisting}
\end{enumerate}